\section*{Executive Summary}
\addcontentsline{toc}{section}{Executive Summary}

\textit{This document is a formal research programme specification with a
fully implemented foundational layer (Cluster~I) and detailed, falsifiable
architectural hypotheses for the remaining layers.}

The \textsc{Apeiron} Framework is a seventeen‑layer architecture for
non‑anthropocentric knowledge genesis.  Its foundational claim is that
an artificial intelligence can construct categories, temporalities,
and causal structures entirely from elementary observables, without
inheriting human labels or objectives.

\paragraph{What has been accomplished.}
\textbf{Cluster~I (Layers~1--4)} is fully specified, implemented, and
empirically validated:
\begin{itemize}
  \item \textbf{Layer~1} defines irreducible observables via
        multi‑axial atomicity, with machine‑checkable certificates
        from SymPy, Z3, Lean~4, and Coq.
  \item \textbf{Layer~2} builds a weighted relational hypergraph from
        pixel co‑activations.
  \item \textbf{Layer~3} identifies functional clusters with higher
        modularity than $k$‑means or spectral clustering across five
        image datasets, and their ablation significantly degrades
        classification accuracy.
  \item \textbf{Layer~4} models endogenous dynamics via stochastic
        differential equations with proven convergence.
  \item A first pilot measurement of \emph{non‑anthropocentricity}
        shows that the discovered clusters are functionally potent
        yet align only weakly with human class labels (NMI as low as
        0.069 on CIFAR‑10).
  \item Full resource profiling confirms that the entire pipeline
        runs in under 8\,minutes on a commodity laptop.
\end{itemize}

\paragraph{What is specified for future work.}
\textbf{Clusters~II--IV (Layers~5--17)} are published as detailed
architectural specifications with falsifiable research questions,
including meta‑learning, ontological creation, reconciliation of
incommensurable ontologies, autopoietic world‑building, responsible
morphogenesis, and absolute integration.  An embedded ethical and
governance framework makes reversibility, auditability, and
multi‑stakeholder deliberation architectural primitives.

\paragraph{Why it matters.}
\textsc{Apeiron} offers a coherent, mathematically rigorous
alternative to the anthropocentric bias of mainstream AI.  It
provides a structured path from formal foundations to planetary‑scale
governance, inviting the research community to collaborate on a
Copernican turn in artificial intelligence.